Bei konstanter Temperatur gilt das Gesetz von Boyle-Mariotte, $p_1V_1 = p_2V_2$. Der Druck sinkt auf
\begin{align}
 p_2 &= p_1 - \dprO\cdot p_1 = \pq\cdot p_1
\end{align}
Damit wird
\begin{align}
 \frac{V_2}{V_1} &= \vqF = \frac{p_1}{\pq\cdot p_1} = \vq
\end{align}
Das Volumen nimmt also zu um
\begin{align}
 \frac{\Delta V}{V_1} &= \dvrF = \dvr \approx \result{\dvrP{0}{2}}
\end{align}
Die Zunahme ist grösser als die Abnahme des Drucks, weil $V \propto 1/p$ ist.
